\section{Method: Trajectory-Anchored Dynamic Selection (TADS)}
\label{sec:method}

We formalise the data-selection problem
(Sec.~\ref{sec:problem-setup}) and the variance issue that motivates
a structural anchor (Sec.~\ref{sec:method-motivation}), introduce
the trajectory-anchor extraction and the composite-reward ranking
rule (Sec.~\ref{sec:method-core}), and finally state the stability
theorem (Sec.~\ref{sec:method-theory}).


\subsection{Problem setup}
\label{sec:problem-setup}

Let $\mathcal{D}=\{x_i\}_{i=1}^N$ be a candidate instruction-tuning
pool, where each $x_i$ is an instruction--response pair tokenised
as a single sequence. Let $f_\theta$ be a base LLM with parameters
$\theta$. Given a selection budget $B \le N$, equivalently a
selection ratio $\rho=B/N$ with $0<\rho\le1$, the goal is to choose
a subset $\mathcal{S}\subset\mathcal{D}$ with $|\mathcal{S}|=B$ so
that fine-tuning $f_\theta$ on $\mathcal{S}$ maximises performance
on evaluation tasks $\{T_1,\ldots,T_M\}$. We work in the
\emph{multi-task} setting, where a single subset $\mathcal{S}$ is
used for all $M$ tasks. \textsc{TADS} treats this as a per-epoch
\emph{dynamic re-ranking} problem: each candidate $x_i$ is assigned
a scalar score $s_i^{(t)}$, and the top-$B$ are selected at each
epoch $t$.

\paragraph{Token-level state dynamics.}
For a hidden layer $l \in \{1,\ldots,L\}$ of $f_\theta$ and an
input sequence $x$ with $K_x$ valid tokens, let
$h_l^{(k)}(x;\theta)\in\mathbb{R}^d$ denote the layer-$l$ hidden
state at the $k$-th token position. The transition between
consecutive token positions is
\begin{equation}
  s_l^{(k)}(x;\theta) := h_l^{(k)}(x;\theta)-h_l^{(k-1)}(x;\theta),
  \qquad k=2,\ldots,K_x .
  \label{eq:transition}
\end{equation}

\paragraph{Contextualisation delta.}
With $h_l^{\mathrm{first}}(x;\theta) := h_l^{(1)}(x;\theta)$ and
$h_l^{\mathrm{last}}(x;\theta) := h_l^{(K_x)}(x;\theta)$,
\begin{equation}
  \Delta h_l(x;\theta) :=
  h_l^{\mathrm{last}}(x;\theta) - h_l^{\mathrm{first}}(x;\theta)
  = \sum_{k=2}^{K_x} s_l^{(k)}(x;\theta).
  \label{eq:delta}
\end{equation}

\paragraph{Sequence-mean activation.}
\begin{equation}
  \bar h_l(x;\theta) :=
  \frac{1}{K_x}\sum_{k=1}^{K_x} h_l^{(k)}(x;\theta).
  \label{eq:meanact}
\end{equation}

\subsection{Why multi-task selection benefits from a structural anchor}
\label{sec:method-motivation}

Composite-reward dynamic selectors~\cite{kung2023activeit,
sun2025dots} score each candidate by combining a difficulty term
$\mathcal{L}_i$ (response-token cross-entropy) and an uncertainty
term $\mathcal{H}_i$ (mean predictive entropy):
\begin{equation}
  \mathcal{R}_i = w\,\mathcal{L}_i + (1-w)\,\mathcal{H}_i,
  \label{eq:composite}
\end{equation}
with the variance-ratio weight
\begin{equation}
  w =
  \frac{\mathrm{Var}_{x \in \mathcal{D}}(\mathcal{L})}
       {\mathrm{Var}_{x \in \mathcal{D}}(\mathcal{L}) +
        \mathrm{Var}_{x \in \mathcal{D}}(\mathcal{H}) + \varepsilon}.
  \label{eq:w}
\end{equation}
We adopt the composite-reward construction unchanged from prior
IT-data selection~\cite{sun2025dots, kung2023activeit}.

\paragraph{The multi-task variance trap.}
By the law of total variance,
\begin{equation}
  \mathrm{Var}_{x \in \mathcal{D}}(\mathcal{L})
  =
  \underbrace{\mathbb{E}_{m}\bigl[\mathrm{Var}_{x \in \mathcal{D}_m}(\mathcal{L})\bigr]}_{\text{within-task}}
  +
  \underbrace{\mathrm{Var}_{m}\bigl[\mathbb{E}_{x \in \mathcal{D}_m}(\mathcal{L})\bigr]}_{\text{between-task}},
  \label{eq:total-variance}
\end{equation}
and analogously for $\mathcal{H}$. In a multi-task pool the
between-task term dominates, so $w$ depends strongly on the task
mix rather than per-sample informativeness and drifts across runs.

\paragraph{A pool-level structural signal complements the scalar reward.}
The trajectory anchor $v_l^{(t)}$ (Sec.~\ref{sec:method-core}) is a
pool-level summary -- the top principal direction of the
contextualisation deltas across a probe subset, not a per-sample
statistic. Its epoch-to-epoch change is uniformly bounded by the
learning rate (Theorem~\ref{thm:stability}), while $w$ has no
comparable bound. \textsc{TADS} composes the task-mix-drifting
scalar reward $\mathcal{R}_i^{(t)}$ (Eq.~\ref{eq:composite}) with
a stability-controlled structural factor: the scalar reward carries
the per-candidate text-level magnitude signal, and the anchor
factor adds a complementary representation-level direction that is
not available from $\mathcal{R}_i^{(t)}$ alone.

\begin{figure}[t]
\centering
\begin{tikzpicture}[scale=0.95, every node/.style={font=\small}]
  \draw[gray!30, very thin, -{Latex[length=2mm]}] (-0.1,0) -- (6.4,0);
  \draw[gray!30, very thin, -{Latex[length=2mm]}] (0,-0.1) -- (0,3.8);
  \node[gray!50, font=\scriptsize] at (6.5,-0.18) {$\mathbb{R}^d$};
  \filldraw[black] (0,0) circle (1.2pt);
  \foreach \x/\y in
    {3.8/1.2, 4.2/1.6, 3.5/1.8, 4.5/2.0, 3.0/1.3, 4.0/2.3, 3.7/1.5, 4.6/1.4, 3.3/1.6}
  {
    \draw[-{Latex[length=1.6mm]}, orange!60, thin, opacity=0.75]
      (0,0) -- (\x,\y);
  }
  \node[orange!80!black, font=\footnotesize, align=left, anchor=north west] at (3.5,0.85)
    {Probe deltas\\
     $\{\Delta h_l(x;\theta_{t-1})\}_{x\in\widetilde{\mathcal{D}}_t}$};
  \draw[-{Latex[length=2.6mm]}, blue!75!black, ultra thick]
    (0,0) -- (5.0,2.0);
  \node[blue!75!black, font=\bfseries] at (5.35,2.25) {$v_l^{(t)}$};
  \draw[-{Latex[length=2.2mm]}, blue!35, dashed, semithick]
    (0,0) -- (5.0,1.65);
  \node[blue!50, font=\scriptsize] at (5.40,1.45) {$v_l^{(t-1)}$};
  \node[gray!60!black, font=\scriptsize, align=left] at (6.4,1.85)
    {Thm.~\ref{thm:stability}:\\[-1pt]
     $\le \tfrac{2C_\Sigma}{\gamma_{\min}}\,\eta_t$};
  \filldraw[red!75!black] (2.2,2.85) circle (2pt);
  \node[red!75!black, font=\footnotesize] at (1.45,3.05)
    {$\bar h_l(x_i;\theta_{t-1})$};
  \draw[red!75!black, dotted, thick] (2.2,2.85) -- (2.88,1.15);
  \filldraw[red!75!black] (2.88,1.15) circle (1.5pt);
  \node[red!75!black, font=\scriptsize, anchor=west]
    at (1.7, 1.2) {$\mathrm{align}_i^{(t)}$};
\end{tikzpicture}
\caption{\textbf{Trajectory anchor at epoch~$t$.}
Orange arrows are the contextualisation deltas
$\Delta h_l(x;\theta_{t-1})$ from a probe subset
$\widetilde{\mathcal{D}}_t$---the \emph{sequence-level} trajectory
inside each sample. The bold blue arrow $v_l^{(t)}$ is their top
principal direction (the anchor); the dashed $v_l^{(t-1)}$ is the
previous epoch's anchor, the \emph{epoch-level} trajectory whose
deviation is bounded by Theorem~\ref{thm:stability}. A candidate's
alignment is the projection of its sequence-mean activation
$\bar h_l(x_i;\theta_{t-1})$ onto $v_l^{(t)}$.}
\label{fig:anchor}
\end{figure}

\subsection{The trajectory anchor and TADS ranking score}
\label{sec:method-core}

The central object of \textsc{TADS} is the \emph{trajectory anchor}
$v_l^{(t)}$: a per-layer capability direction extracted online
from the model's own hidden-state trajectory at each refresh step.
We first describe how the anchor is extracted
(Sec.~\ref{sec:method-core}, ``Probe subset and trajectory
anchor'') and how a candidate's alignment to it is measured
(``Candidate alignment''); these are the novel ingredients of the
selection rule. We then describe a supporting deterministic
transform of the composite reward into a bounded composite reward
(``Composite reward''), which serves as the carrier signal that
the anchor factor modulates, and combine the two into the final
score (``Selection score''), controlled by a single hyperparameter
$\lambda \ge 0$.

\paragraph{Probe subset and trajectory anchor.}
At the start of epoch~$t$, we draw a probe subset
$\widetilde{\mathcal{D}}_t \subset \mathcal{D}$ uniformly at random
with $|\widetilde{\mathcal{D}}_t| = 1024$, re-sampled each epoch.
For each layer $l$, we collect
$\{\Delta h_l(x;\theta_{t-1})\}_{x \in \widetilde{\mathcal{D}}_t}$
and form the empirical mean and centred covariance:
\begin{align}
  \overline{\Delta h}_l^{(t)}
   &:= \frac{1}{|\widetilde{\mathcal{D}}_t|}
        \sum_{x \in \widetilde{\mathcal{D}}_t}
        \Delta h_l(x;\theta_{t-1}),
   \label{eq:mu} \\[2pt]
  \Sigma_l^{(t)}
   &:= \frac{1}{|\widetilde{\mathcal{D}}_t|}
        \sum_{x \in \widetilde{\mathcal{D}}_t}
        \bigl(\Delta h_l - \overline{\Delta h}_l^{(t)}\bigr)
        \bigl(\Delta h_l - \overline{\Delta h}_l^{(t)}\bigr)^{\!\top}.
   \label{eq:sigma}
\end{align}
The trajectory anchor at epoch~$t$ for layer $l$ is the top unit
eigenvector of $\Sigma_l^{(t)}$:
\begin{equation}
  v_l^{(t)} = \arg\max_{\|v\| = 1} v^{\!\top}\Sigma_l^{(t)}v ,
  \label{eq:anchor-def}
\end{equation}
sign-calibrated so that
$\langle v_l^{(t)}, \overline{\Delta h}_l^{(t)} \rangle > 0$.
The anchor depends only on $\theta_{t-1}$ and the candidate pool;
it requires no in-domain seed samples, labels, or auxiliary model.

\paragraph{Candidate alignment.}
The per-layer alignment is
$\langle \bar h_l(x_i;\theta_{t-1}), v_l^{(t)} \rangle$. Aggregating
across layers and normalising,
\begin{equation}
  \mathrm{align}_i^{(t)} :=
  \frac{1}{L}\sum_{l=1}^{L}
  \bigl\langle \bar h_l(x_i;\theta_{t-1}), v_l^{(t)} \bigr\rangle,
  \label{eq:align-raw}
\end{equation}
\begin{equation}
  \widetilde{\mathrm{align}}_i^{(t)} :=
  \frac{\mathrm{align}_i^{(t)} - \min_j \mathrm{align}_j^{(t)}}
       {\max_j \mathrm{align}_j^{(t)} - \min_j \mathrm{align}_j^{(t)}}
  \in [0,1].
  \label{eq:align}
\end{equation}

\paragraph{Selection score.}
\textsc{TADS} forms the final per-candidate ranking score by
multiplying the composite reward by the trajectory-anchor factor:
\begin{equation}
\boxed{
s_i^{(t)}
=
\underbrace{\mathcal{R}_i^{(t)}}_{\substack{\text{composite-reward factor}\\[1pt]
\text{sample-level signal}}}
\cdot
\underbrace{\bigl(1+\lambda \cdot \widetilde{\mathrm{align}}_i^{(t)}\bigr)}_{\substack{\text{anchor factor} \in [1, 1+\lambda]\\[1pt]
\text{stability-controlled pool-level signal}}}
}
\label{eq:score}
\end{equation}
with hyperparameter $\lambda \ge 0$. The composite reward
$\mathcal{R}_i^{(t)}$ is used as-is from prior dynamic
selectors~\cite{kung2023activeit, sun2025dots}: no learned
transform, no learned policy, no z-score or sigmoid is applied
inside the ranking rule. The multiplicative form carries three
useful properties. \emph{(i) Clean ablation:} $\lambda{=}0$
recovers the composite-reward base ranking $\mathcal{R}_i^{(t)}$
exactly; any gain at $\lambda>0$ is attributable to the
trajectory anchor. \emph{(ii) Signal separation:} the
composite-reward factor carries text-level magnitude information
($\mathcal{L}, \mathcal{H}$) and inherits any task-mix drift of
Eq.~\ref{eq:total-variance}; the anchor factor is a bounded
representation-level pool-level PCA statistic. The two factors
operate on different channels (text vs. hidden state) and the
multiplicative form composes them at the relative scale where
each is informative, which an additive composition would forfeit.
\emph{(iii) Stability transfer:} the anchor factor is bounded in
$[1, 1+\lambda]$ and its trajectory anchor is stable by
Theorem~\ref{thm:stability}; the anchor-induced change in
$s_i^{(t)}$ is bounded by $O(\lambda\eta_t)$
(App.~\ref{app:score-stability}).


\paragraph{Implementation note: dataset-level $w$.}
The variance-ratio weight $w$ in Eq.~\ref{eq:w} is computed at the
dataset level. A naive mini-batch implementation degenerates to
$w \equiv 0$ when $|\mathcal{B}|=1$, silently collapsing the
difficulty signal. \textsc{TADS} accumulates $\{\mathcal{L}_i\}$
and $\{\mathcal{H}_i\}$ across all batches and computes $w$ once
per epoch, then forms $\mathcal{R}_i^{(t)}$ via
Eq.~\ref{eq:composite}. No additional forward passes are required
beyond those used to collect $\{\mathcal{L}_i, \mathcal{H}_i\}$.

\begin{algorithm}[t]
\caption{Trajectory-Anchored Dynamic Selection (TADS)}
\label{alg:tads}
\begin{algorithmic}[1]
\REQUIRE Pool $\mathcal{D}$, base model $f_{\theta_0}$, budget $B$,
         layers $\{1,\ldots,L\}$, weight $\lambda$, epochs $T$, probe
         size $|\widetilde{\mathcal{D}}|$
\FOR{$t = 1, \dots, T$}
  \STATE Sample probe set $\widetilde{\mathcal{D}}_t$ uniformly at random
  \FOR{$l = 1, \dots, L$}
    \STATE Compute deltas
           $\{\Delta h_l(x; \theta_{t-1})\}_{x \in \widetilde{\mathcal{D}}_t}$
           and covariance $\Sigma_l^{(t)}$
           \hfill (Eqs.~\ref{eq:delta}--\ref{eq:sigma})
    \STATE $v_l^{(t)} \leftarrow$ top-1 eigenvector of
           $\Sigma_l^{(t)}$, sign-calibrated
  \ENDFOR
  \STATE For all $i$: aggregate $\widetilde{\mathrm{align}}_i^{(t)}$
         \hfill (Eq.~\ref{eq:align})
  \STATE Forward $\theta_{t-1}$ over $\mathcal{D}$ to collect
         $\{\mathcal{L}_i, \mathcal{H}_i\}_{i=1}^N$
  \STATE Compute pool-level $w^{(t)}$ and composite reward
         $\mathcal{R}_i^{(t)} = w^{(t)}\mathcal{L}_i + (1-w^{(t)})\mathcal{H}_i$
         \hfill (Eqs.~\ref{eq:w},~\ref{eq:composite})
  \STATE Score
         $s_i^{(t)} \leftarrow
          \mathcal{R}_i^{(t)} \cdot
          (1 + \lambda \cdot \widetilde{\mathrm{align}}_i^{(t)})$
         \hfill (Eq.~\ref{eq:score})
  \STATE $\mathcal{S}_t \leftarrow$ top-$B$ samples by $s_i^{(t)}$;
         fine-tune $\theta_{t-1} \to \theta_t$ on $\mathcal{S}_t$
\ENDFOR
\RETURN $\theta_T$
\end{algorithmic}
\end{algorithm}

\subsection{Stability of the trajectory anchor}
\label{sec:method-theory}

\begin{theorem}[Anchor stability]
\label{thm:stability}
Fix a layer $l$ and write $\Sigma^{(t)} := \Sigma_l^{(t)}$. Assume
\emph{(1)} bounded covariance change,
$\|\Sigma^{(t+1)}-\Sigma^{(t)}\|_F \le C_\Sigma \eta_t$;
\emph{(2)} a positive eigenvalue gap,
$\gamma^{(t)} := \lambda_1^{(t)}-\lambda_2^{(t)}
\ge \gamma_{\min}>0$; and \emph{(3)} consistent sign calibration.
Then the sign-calibrated trajectory anchor satisfies
\begin{equation}
  \boxed{\bigl\|v_l^{(t+1)}-v_l^{(t)}\bigr\|
  \le \frac{2C_\Sigma}{\gamma_{\min}}\eta_t}
  \label{eq:stability-bound}
\end{equation}
and, by summing over epochs,
$\|v_l^{(t+m)}-v_l^{(t)}\| \le (2C_\Sigma/\gamma_{\min})
\sum_{s=t}^{t+m-1}\eta_s.$
\end{theorem}

The proof applies the Davis--Kahan $\sin\Theta$
theorem~\cite{daviskahan}; details in
App.~\ref{app:proof-stability}. The corollary applies to any
uniformly bounded base ranking score, so the stability statement
for the \textsc{TADS} criterion is not tied to the composite-reward
form of Eq.~\ref{eq:composite}: any bounded per-candidate signal
on the left of Eq.~\ref{eq:score} would inherit the same
$O(\lambda\eta_t)$ anchor-induced perturbation. Empirical
verification across 1{,}496 layer-refresh measurements
(App.~\ref{app:thm-verification}) confirms the bound is non-vacuous.
